\documentclass[stu]{apa7}

\usepackage[spanish]{babel}
\usepackage[utf8]{inputenc}
\usepackage{hyperref}
\usepackage{enumitem}
\usepackage{geometry}
\usepackage{times}
\usepackage{array}
\usepackage{longtable}
\usepackage{booktabs}
\usepackage{colortbl}
\usepackage{xcolor}
\usepackage{tabularx}

\geometry{margin=1in}

\title{Casos de Uso --- Sistema de información web contable sobre recibos de retención y movimientos de caja de la empresa minera MINCH SRL.}
\author{...}
\affiliation{...}
\course{...}
\professor{...}
\duedate{...}

\setlength{\parindent}{0.5in}

\newcolumntype{L}[1]{>{\raggedright\arraybackslash}p{#1}}
\newcolumntype{C}[1]{>{\centering\arraybackslash}p{#1}}

\begin{document}

\maketitle

\section{Perfiles de usuario}

\begin{longtable}{|L{3.5cm}|L{6cm}|L{3cm}|L{3cm}|}
\hline
\rowcolor[gray]{0.9}
\textbf{Actor} & \textbf{Responsabilidad clave (Dominio)} & \textbf{Nivel de acceso} & \textbf{Rol en el sistema} \\
\hline
\endhead
Administrador del Sistema & Gestión de infraestructura, seguridad, creación de usuarios, asignación de roles y permisos. & Total (Root) & Admin \\
\hline
Contador General & Gestión de retenciones (cálculo automático), conciliación de libro de bancos, estados de cuenta por cliente y proveedor, supervisión de caja. & Estratégico & Contador \\
\hline
Auxiliar Contable & Registro diario de movimientos de caja, emisión de recibos de retención, ingreso de movimientos bancarios. & Operativo & Auxiliar \\
\hline
Gerente / Administrador & Visualización de dashboard ejecutivo, generación de reportes PDF y Excel, toma de decisiones financieras. & Estratégico & Gerente \\
\hline
Proveedor & Destinatario de recibos de retención y estados de cuenta (información consultada y entregada por el personal). & Externo & --- \\
\hline
Cliente & Destinatario de estados de cuenta. & Externo & --- \\
\hline
\end{longtable}

\newpage

\section{Diagrama de casos de uso}

\subsection{Convenciones}

Para la elaboración del diagrama de casos de uso, se utilizan las siguientes relaciones:

\begin{itemize}
    \item \textbf{Asociación:} Línea continua entre actor y caso de uso. Indica que el actor participa en ese caso de uso.
    \item \textbf{«include»:} Relación obligatoria. Un caso de uso base incluye siempre la funcionalidad del caso incluido. Se dibuja como flecha discontinua con estereotipo «include» desde el caso base hacia el caso incluido.
    \item \textbf{«extend»:} Relación opcional. Un caso de uso extiende al caso base bajo ciertas condiciones. Se dibuja como flecha discontinua con estereotipo «extend» desde el caso extendido hacia el caso base.
    \item \textbf{Generalización:} Flecha con punta vacía de un actor especializado hacia un actor general.
\end{itemize}

\subsection{Módulo de autenticación}

\begin{longtable}{|L{3cm}|L{4cm}|L{4cm}|L{4cm}|}
\hline
\rowcolor[gray]{0.9}
\textbf{Actor} & \textbf{Caso de uso} & \textbf{Tipo de interacción} & \textbf{Relacionado con} \\
\hline
\endhead
Todos los actores & Iniciar sesión & Asociación directa & --- \\
\hline
\end{longtable}

\textbf{Instrucciones para el diagrama:}
Dibujar un rectángulo etiquetado \textbf{«Sistema MINCH SRL»}. Fuera del rectángulo, colocar los actores: Administrador del Sistema, Contador General, Auxiliar Contable, Gerente. Dentro del rectángulo, dibujar un óvalo etiquetado \textbf{«Iniciar sesión»}. Conectar cada actor a este óvalo con una línea continua.

\subsection{Módulo de retenciones}

\begin{longtable}{|L{3cm}|L{5cm}|L{4cm}|L{3cm}|}
\hline
\rowcolor[gray]{0.9}
\textbf{Actor} & \textbf{Caso de uso} & \textbf{Tipo de interacción} & \textbf{Relacionado con} \\
\hline
\endhead
Contador General & Gestionar retenciones & Asociación directa & --- \\
\hline
--- & Calcular descuentos automáticamente & «include» & Gestionar retenciones \\
\hline
Contador General & Generar recibo de retención PDF & «extend» & Gestionar retenciones \\
\hline
Contador General & Exportar retenciones a Excel & «extend» & Gestionar retenciones \\
\hline
\end{longtable}

\textbf{Instrucciones para el diagrama:}
Dentro del sistema, dibujar un óvalo \textbf{«Gestionar retenciones»}. Conectar al actor Contador General con línea continua. Dibujar otro óvalo \textbf{«Calcular descuentos automáticamente»} y conectar con flecha discontinua «include» desde \textbf{Gestionar retenciones} hacia \textbf{Calcular descuentos}. Dibujar óvalos \textbf{«Generar recibo de retención PDF»} y \textbf{«Exportar retenciones a Excel»}, cada uno con flecha discontinua «extend» desde ellos hacia \textbf{Gestionar retenciones}.

\subsection{Módulo de libro de bancos}

\begin{longtable}{|L{3cm}|L{5cm}|L{4cm}|L{3cm}|}
\hline
\rowcolor[gray]{0.9}
\textbf{Actor} & \textbf{Caso de uso} & \textbf{Tipo de interacción} & \textbf{Relacionado con} \\
\hline
\endhead
Auxiliar Contable & Gestionar movimientos bancarios & Asociación directa & --- \\
\hline
--- & Numerar transacción automáticamente & «include» & Gestionar movimientos bancarios \\
\hline
Auxiliar Contable & Generar recibo de banco PDF & «extend» & Gestionar movimientos bancarios \\
\hline
Contador General & Gestionar cuentas bancarias & Asociación directa & --- \\
\hline
\end{longtable}

\textbf{Instrucciones para el diagrama:}
Dibujar óvalos \textbf{«Gestionar movimientos bancarios»} (conectado a Auxiliar Contable), \textbf{«Numerar transacción automáticamente»} (con «include» desde Gestionar movimientos bancarios), \textbf{«Generar recibo de banco PDF»} (con «extend» hacia Gestionar movimientos bancarios), y \textbf{«Gestionar cuentas bancarias»} (conectado a Contador General).

\subsection{Módulo de caja}

\begin{longtable}{|L{3cm}|L{5cm}|L{4cm}|L{3cm}|}
\hline
\rowcolor[gray]{0.9}
\textbf{Actor} & \textbf{Caso de uso} & \textbf{Tipo de interacción} & \textbf{Relacionado con} \\
\hline
\endhead
Auxiliar Contable & Gestionar movimientos de caja & Asociación directa & --- \\
\hline
--- & Numerar movimiento de caja automáticamente & «include» & Gestionar movimientos de caja \\
\hline
Auxiliar Contable & Generar recibo de caja PDF & «extend» & Gestionar movimientos de caja \\
\hline
\end{longtable}

\textbf{Instrucciones para el diagrama:}
Dibujar óvalo \textbf{«Gestionar movimientos de caja»} (conectado a Auxiliar Contable), \textbf{«Numerar movimiento de caja automáticamente»} (con «include»), \textbf{«Generar recibo de caja PDF»} (con «extend»).

\subsection{Módulo de estados de cuenta}

\begin{longtable}{|L{3cm}|L{5cm}|L{4cm}|L{3cm}|}
\hline
\rowcolor[gray]{0.9}
\textbf{Actor} & \textbf{Caso de uso} & \textbf{Tipo de interacción} & \textbf{Relacionado con} \\
\hline
\endhead
Contador General & Consultar estado de cuenta & Asociación directa & --- \\
\hline
Contador General & Generar PDF estado de cuenta & «extend» & Consultar estado de cuenta \\
\hline
\end{longtable}

\textbf{Instrucciones para el diagrama:}
Dibujar óvalo \textbf{«Consultar estado de cuenta»} (conectado a Contador General), \textbf{«Generar PDF estado de cuenta»} (con «extend» hacia Consultar estado de cuenta). Los actores externos Proveedor y Cliente se colocan fuera del sistema sin conexión directa a casos de uso (reciben la información a trav\'es del Contador).

\subsection{Módulo de reportes}

\begin{longtable}{|L{3cm}|L{5cm}|L{4cm}|L{3cm}|}
\hline
\rowcolor[gray]{0.9}
\textbf{Actor} & \textbf{Caso de uso} & \textbf{Tipo de interacción} & \textbf{Relacionado con} \\
\hline
\endhead
Contador General & Generar reportes & Asociación directa & --- \\
\hline
Gerente & Visualizar dashboard & Asociación directa & --- \\
\hline
Contador General / Gerente & Exportar reporte a Excel & «extend» & Generar reportes \\
\hline
\end{longtable}

\textbf{Instrucciones para el diagrama:}
Dibujar óvalo \textbf{«Generar reportes»} (conectado a Contador General), \textbf{«Visualizar dashboard»} (conectado a Gerente), \textbf{«Exportar reporte a Excel»} (con «extend» hacia Generar reportes).

\subsection{Módulo de administración}

\begin{longtable}{|L{3cm}|L{5cm}|L{4cm}|L{3cm}|}
\hline
\rowcolor[gray]{0.9}
\textbf{Actor} & \textbf{Caso de uso} & \textbf{Tipo de interacción} & \textbf{Relacionado con} \\
\hline
\endhead
Administrador del Sistema & Gestionar usuarios & Asociación directa & --- \\
\hline
Administrador del Sistema & Gestionar roles y permisos & Asociación directa & --- \\
\hline
Contador General & Gestionar proveedores & Asociación directa & --- \\
\hline
Contador General & Gestionar clientes & Asociación directa & --- \\
\hline
Contador General & Gestionar impuestos/tasas & Asociación directa & --- \\
\hline
\end{longtable}

\textbf{Instrucciones para el diagrama:}
Dibujar los óvalos correspondientes en el área de administración del sistema. Conectar cada uno al actor respectivo. No existen relaciones «include» ni «extend» entre estos casos de uso, son independientes.

\newpage

\section{Especificación de casos de uso detallados}

\subsection{Módulo de autenticación}

\subsubsection{Caso de Uso CU-01: Iniciar sesión}

\begin{tabularx}{\textwidth}{|l|X|}
\hline
\rowcolor[gray]{0.9}
\textbf{Campo} & \textbf{Valor} \\
\hline
Caso de Uso & Iniciar sesión \\
\hline
Actor & Todos los actores (Administrador, Contador, Auxiliar, Gerente) \\
\hline
Descripción & El usuario ingresa al sistema proporcionando sus credenciales (email y contraseña). El sistema verifica la identidad y redirige al dashboard según el rol y los permisos asignados. \\
\hline
\end{tabularx}

\vspace{0.3cm}

\textbf{Flujo de eventos básico}

\begin{tabularx}{\textwidth}{|c|X|}
\hline
& \textbf{Usuario} \\
\hline
1 & Ingresa a la URL del sistema. \\
\hline
2 & Introduce su dirección de email y contraseña. \\
\hline
3 & Presiona el botón \textbf{Iniciar sesión}. \\
\hline
& \textbf{Sistema} \\
\hline
4 & Valida que el email y la contraseña sean correctos. \\
\hline
5 & Verifica que el usuario tenga un rol activo con permisos asignados. \\
\hline
6 & Redirige al dashboard principal del sistema. \\
\hline
\end{tabularx}

\vspace{0.3cm}

\textbf{Flujo alterno}

\begin{tabularx}{\textwidth}{|c|X|}
\hline
& \textbf{Sistema} \\
\hline
4.1 & Si el email o la contraseña son incorrectos, muestra el mensaje: \textit{«Estas credenciales no coinciden con nuestros registros.»} \\
\hline
4.2 & Si la cuenta está bloqueada o deshabilitada, muestra: \textit{«Cuenta deshabilitada. Contacte al administrador.»} \\
\hline
4.3 & Si el usuario no tiene ningún rol asignado, muestra: \textit{«Acceso denegado. No tiene permisos asignados.»} \\
\hline
\end{tabularx}

\vspace{0.3cm}

\textbf{Precondiciones:} El usuario debe estar registrado en el sistema con un email válido y una contraseña encriptada. Debe tener al menos un rol con permisos asignados.

\textbf{Postcondición:} El usuario accede al sistema con una sesión activa. Se registra la fecha y hora del último acceso.

\newpage

\subsection{Módulo de retenciones}

\subsubsection{Caso de Uso CU-02: Gestionar retenciones}

\begin{tabularx}{\textwidth}{|l|X|}
\hline
\rowcolor[gray]{0.9}
\textbf{Campo} & \textbf{Valor} \\
\hline
Caso de Uso & Gestionar retenciones \\
\hline
Actor & Contador General \\
\hline
Descripción & El Contador registra, modifica, consulta o elimina retenciones. Cada retención corresponde a un proveedor sin NIT al que se le retiene un porcentaje según el tipo de bien o servicio. El sistema calcula automáticamente los descuentos aplicables (RC-IVA, IUE, IT) y genera el código correlativo mensual. \\
\hline
\end{tabularx}

\vspace{0.3cm}

\textbf{Flujo de eventos básico}

\begin{tabularx}{\textwidth}{|c|X|}
\hline
& \textbf{Usuario} \\
\hline
1 & Ingresa al módulo de retenciones. \\
\hline
2 & Presiona el botón \textbf{Nueva retención}. \\
\hline
3 & Selecciona el tipo de retención (Servicios o Bienes). \\
\hline
4 & Busca y selecciona un proveedor por CI o nombre. \\
\hline
5 & Ingresa el monto base de la retención. \\
\hline
6 & Ingresa la fecha, descripción y resumen. \\
\hline
7 & Presiona el botón \textbf{Guardar}. \\
\hline
& \textbf{Sistema} \\
\hline
8 & Ejecuta \textbf{Calcular descuentos automáticamente} (CU-03) para obtener los descuentos aplicables según las tasas configuradas para el tipo de retención. \\
\hline
9 & Genera el código correlativo automático para el mes en curso. \\
\hline
10 & Guarda la retención, los descuentos calculados y el código en la base de datos. \\
\hline
11 & Muestra un mensaje de confirmación: \textit{«Retención registrada correctamente.»} \\
\hline
\end{tabularx}

\vspace{0.3cm}

\textbf{Flujo alterno}

\begin{tabularx}{\textwidth}{|c|X|}
\hline
& \textbf{Sistema} \\
\hline
7.1 & Si el monto ingresado está fuera del rango permitido (10.00 -- 999,999,999.99), muestra: \textit{«El monto debe estar entre 10.00 y 999,999,999.99.»} \\
\hline
7.2 & Si no selecciona un proveedor, muestra: \textit{«Debe seleccionar un proveedor.»} \\
\hline
7.3 & Si el proveedor no existe, el sistema permite registrarlo automáticamente mediante la búsqueda por CI. \\
\hline
\end{tabularx}

\vspace{0.3cm}

\textbf{Precondiciones:} El Contador debe haber iniciado sesión. Deben existir impuestos/tasas configurados (RC-IVA, IUE, IT). Debe existir al menos un proveedor registrado.

\textbf{Postcondición:} La retención queda registrada con su código correlativo, los descuentos calculados y el estado pendiente. Se puede generar el recibo PDF (CU-04).

\newpage

\subsubsection{Caso de Uso CU-03: Calcular descuentos automáticamente}

\begin{tabularx}{\textwidth}{|l|X|}
\hline
\rowcolor[gray]{0.9}
\textbf{Campo} & \textbf{Valor} \\
\hline
Caso de Uso & Calcular descuentos automáticamente \\
\hline
Actor & --- (Incluido en CU-02) \\
\hline
Descripción & Este caso de uso es ejecutado internamente por \textbf{Gestionar retenciones} (CU-02). Calcula los descuentos aplicables a una retención según las tasas impositivas configuradas para el tipo de bien o servicio. \\
\hline
\end{tabularx}

\vspace{0.3cm}

\textbf{Flujo de eventos básico}

\begin{tabularx}{\textwidth}{|c|X|}
\hline
& \textbf{Sistema} \\
\hline
1 & Recibe el tipo de retención (S = Servicios o G = Bienes) y el monto base. \\
\hline
2 & Consulta las tasas de impuestos aplicables: para Servicios (RC-IVA 13\% + IT 3\%), para Bienes (IUE 5\% + IT 3\%). \\
\hline
3 & Calcula el total base: \textit{total = monto / (1 - suma\_tasas / 100)}. \\
\hline
4 & Calcula cada descuento individual: \textit{descuento = total * tasa / 100}. \\
\hline
5 & Retorna los descuentos calculados para ser almacenados junto con la retención. \\
\hline
\end{tabularx}

\vspace{0.3cm}

\textbf{Precondiciones:} Las tasas de impuestos deben estar configuradas en el sistema. El tipo de retención (S/G) debe ser válido.

\textbf{Postcondición:} Se retornan los valores calculados de descuentos. No se persiste nada por sí mismo; los resultados son utilizados por CU-02.

\newpage

\subsubsection{Caso de Uso CU-04: Generar recibo de retención PDF}

\begin{tabularx}{\textwidth}{|l|X|}
\hline
\rowcolor[gray]{0.9}
\textbf{Campo} & \textbf{Valor} \\
\hline
Caso de Uso & Generar recibo de retención PDF \\
\hline
Actor & Contador General \\
\hline
Descripción & El Contador genera un documento PDF con el detalle de una retención específica, incluyendo datos del proveedor, montos, descuentos calculados y el código de retención. El PDF se visualiza en el navegador y puede imprimirse o descargarse. \\
\hline
\end{tabularx}

\vspace{0.3cm}

\textbf{Flujo de eventos básico}

\begin{tabularx}{\textwidth}{|c|X|}
\hline
& \textbf{Usuario} \\
\hline
1 & Se encuentra en la lista de retenciones (CU-02). \\
\hline
2 & Presiona el botón \textbf{PDF} en la fila de una retención específica. \\
\hline
& \textbf{Sistema} \\
\hline
3 & Recupera los datos completos de la retención, incluyendo proveedor, descuentos y código. \\
\hline
4 & Renderiza la vista PDF con el formato establecido (tamaño A4, márgenes 15mm). \\
\hline
5 & Genera el documento PDF utilizando la librería mPDF. \\
\hline
6 & Muestra el PDF en el navegador para su visualización, impresión o descarga. \\
\hline
\end{tabularx}

\vspace{0.3cm}

\textbf{Precondiciones:} La retención debe existir y estar registrada con todos sus datos.

\textbf{Postcondición:} El PDF se muestra en pantalla. No se modifica ningún dato en la base de datos.

\newpage

\subsection{Módulo de libro de bancos}

\subsubsection{Caso de Uso CU-05: Gestionar movimientos bancarios}

\begin{tabularx}{\textwidth}{|l|X|}
\hline
\rowcolor[gray]{0.9}
\textbf{Campo} & \textbf{Valor} \\
\hline
Caso de Uso & Gestionar movimientos bancarios \\
\hline
Actor & Auxiliar Contable \\
\hline
Descripción & El Auxiliar registra, modifica, consulta o elimina movimientos bancarios (ingresos y egresos) asociados a una cuenta bancaria. El sistema asigna automáticamente el número de transacción correlativo por año fiscal (octubre -- septiembre) y tipo de movimiento. \\
\hline
\end{tabularx}

\vspace{0.3cm}

\textbf{Flujo de eventos básico}

\begin{tabularx}{\textwidth}{|c|X|}
\hline
& \textbf{Usuario} \\
\hline
1 & Ingresa al módulo de libro de bancos y selecciona una cuenta bancaria y un mes. \\
\hline
2 & Presiona el botón \textbf{Nuevo movimiento}. \\
\hline
3 & Selecciona el tipo de movimiento (D = Débito / C = Crédito). \\
\hline
4 & Busca y selecciona un proveedor o persona asociada. \\
\hline
5 & Ingresa el monto, la fecha, la descripción, el tipo de pago (T = Transferencia / CH = Cheque) y el número de cheque si corresponde. \\
\hline
6 & Presiona \textbf{Guardar}. \\
\hline
& \textbf{Sistema} \\
\hline
7 & Ejecuta \textbf{Numerar transacción automáticamente} (CU-06). \\
\hline
8 & Calcula el nuevo saldo de la cuenta sumando o restando el monto al saldo anterior. \\
\hline
9 & Guarda el movimiento bancario y la transacción en la base de datos. \\
\hline
10 & Recalcula los saldos de meses siguientes si es necesario. \\
\hline
11 & Muestra mensaje de confirmación. \\
\hline
\end{tabularx}

\vspace{0.3cm}

\textbf{Flujo alterno}

\begin{tabularx}{\textwidth}{|c|X|}
\hline
& \textbf{Sistema} \\
\hline
6.1 & Si el tipo de pago es CH (Cheque) y no ingresa el número de cheque, muestra: \textit{«Debe ingresar el número de cheque.»} \\
\hline
6.2 & Si no selecciona un proveedor, permite registrar uno nuevo mediante la búsqueda por CI. \\
\hline
\end{tabularx}

\vspace{0.3cm}

\textbf{Precondiciones:} Debe existir al menos una cuenta bancaria registrada. El usuario debe estar autenticado.

\textbf{Postcondición:} El movimiento queda registrado con su número de transacción. El saldo de la cuenta se actualiza. Los saldos de meses siguientes se recalculan automáticamente.

\newpage

\subsubsection{Caso de Uso CU-06: Numerar transacción automáticamente}

\begin{tabularx}{\textwidth}{|l|X|}
\hline
\rowcolor[gray]{0.9}
\textbf{Campo} & \textbf{Valor} \\
\hline
Caso de Uso & Numerar transacción automáticamente \\
\hline
Actor & --- (Incluido en CU-05) \\
\hline
Descripción & Ejecutado internamente por \textbf{Gestionar movimientos bancarios} (CU-05). Asigna un número correlativo a cada transacción bancaria según el año fiscal (octubre a septiembre) y el tipo de movimiento (débito o crédito), utilizando bloqueo de fila para evitar duplicados en entornos concurrentes. \\
\hline
\end{tabularx}

\vspace{0.3cm}

\textbf{Flujo de eventos básico}

\begin{tabularx}{\textwidth}{|c|X|}
\hline
& \textbf{Sistema} \\
\hline
1 & Recibe el tipo de movimiento (D/C) y la fecha del movimiento. \\
\hline
2 & Determina el año fiscal: si el mes es octubre o posterior, el año fiscal inicia en octubre; si es septiembre o anterior, inicia en octubre del año anterior. \\
\hline
3 & Bloquea la tabla de transacciones (\texttt{lockForUpdate}) para evitar concurrencia. \\
\hline
4 & Consulta el número máximo existente para el mismo tipo y año fiscal. \\
\hline
5 & Asigna el siguiente número correlativo. \\
\hline
6 & Retorna el número asignado. \\
\hline
\end{tabularx}

\vspace{0.3cm}

\textbf{Precondiciones:} La tabla de transacciones debe estar accesible. La fecha del movimiento debe ser válida.

\textbf{Postcondición:} Se retorna un número entero secuencial único para el tipo y año fiscal.

\newpage

\subsection{Módulo de caja}

\subsubsection{Caso de Uso CU-08: Gestionar movimientos de caja}

\begin{tabularx}{\textwidth}{|l|X|}
\hline
\rowcolor[gray]{0.9}
\textbf{Campo} & \textbf{Valor} \\
\hline
Caso de Uso & Gestionar movimientos de caja \\
\hline
Actor & Auxiliar Contable \\
\hline
Descripción & El Auxiliar registra, modifica, consulta o elimina movimientos de caja (ingresos y egresos en efectivo). El sistema asigna automáticamente el número de caja correlativo por mes calendario y tipo de movimiento. \\
\hline
\end{tabularx}

\vspace{0.3cm}

\textbf{Flujo de eventos básico}

\begin{tabularx}{\textwidth}{|c|X|}
\hline
& \textbf{Usuario} \\
\hline
1 & Ingresa al módulo de caja. \\
\hline
2 & Presiona \textbf{Nuevo movimiento}. \\
\hline
3 & Selecciona el tipo (D = Débito / C = Crédito). \\
\hline
4 & Busca y selecciona la persona o proveedor. \\
\hline
5 & Ingresa monto, fecha, descripción. \\
\hline
6 & Presiona \textbf{Guardar}. \\
\hline
& \textbf{Sistema} \\
\hline
7 & Ejecuta \textbf{Numerar movimiento de caja automáticamente} (CU-09). \\
\hline
8 & Si es el primer movimiento del mes, calcula el saldo anterior automáticamente. \\
\hline
9 & Guarda el movimiento y la caja en la base de datos. \\
\hline
10 & Recalcula saldos de meses siguientes. \\
\hline
11 & Muestra mensaje de confirmación. \\
\hline
\end{tabularx}

\vspace{0.3cm}

\textbf{Precondiciones:} El usuario debe estar autenticado.

\textbf{Postcondición:} El movimiento queda registrado con su número de caja. Los saldos se actualizan.

\newpage

\subsection{Módulo de estados de cuenta}

\subsubsection{Caso de Uso CU-11: Consultar estado de cuenta}

\begin{tabularx}{\textwidth}{|l|X|}
\hline
\rowcolor[gray]{0.9}
\textbf{Campo} & \textbf{Valor} \\
\hline
Caso de Uso & Consultar estado de cuenta \\
\hline
Actor & Contador General \\
\hline
Descripción & El Contador consulta el estado de cuenta de un proveedor o cliente. El sistema muestra el detalle de todos los movimientos (débitos y créditos) ordenados por fecha, con el saldo corriente calculado en cada fila, más los totales de débitos, créditos y saldo final. \\
\hline
\end{tabularx}

\vspace{0.3cm}

\textbf{Flujo de eventos básico}

\begin{tabularx}{\textwidth}{|c|X|}
\hline
& \textbf{Usuario} \\
\hline
1 & Ingresa al módulo de estados de cuenta. \\
\hline
2 & Selecciona la pestaña \textbf{Proveedores} o \textbf{Clientes}. \\
\hline
3 & Presiona \textbf{Ver estado de cuenta} en la fila del titular deseado. \\
\hline
& \textbf{Sistema} \\
\hline
4 & Consulta todos los movimientos asociados a la persona (proveedor o cliente). \\
\hline
5 & Calcula el saldo corriente para cada movimiento usando una función de ventana. \\
\hline
6 & Calcula los totales: suma de débitos, suma de créditos, saldo final. \\
\hline
7 & Muestra la tabla de movimientos paginada con los totales. \\
\hline
\end{tabularx}

\vspace{0.3cm}

\textbf{Precondiciones:} Debe existir al menos un proveedor o cliente con movimientos registrados.

\textbf{Postcondición:} Se visualiza el estado de cuenta. No se modifica ningún dato.

\newpage

\subsection{Módulo de reportes y dashboard}

\subsubsection{Caso de Uso CU-13: Generar reportes}

\begin{tabularx}{\textwidth}{|l|X|}
\hline
\rowcolor[gray]{0.9}
\textbf{Campo} & \textbf{Valor} \\
\hline
Caso de Uso & Generar reportes \\
\hline
Actor & Contador General \\
\hline
Descripción & El Contador genera reportes de retenciones, libro de bancos o movimientos de caja aplicando filtros por rango de fechas, tipo y cuenta. El sistema permite visualizar los datos en pantalla y exportarlos a PDF o Excel. \\
\hline
\end{tabularx}

\vspace{0.3cm}

\textbf{Flujo de eventos básico}

\begin{tabularx}{\textwidth}{|c|X|}
\hline
& \textbf{Usuario} \\
\hline
1 & Ingresa al módulo de reportes. \\
\hline
2 & Selecciona el tipo de reporte (Retenciones, Caja, Banco). \\
\hline
3 & Configura los filtros: rango de fechas, tipo de movimiento, cuenta (si aplica). \\
\hline
4 & Presiona \textbf{Generar}. \\
\hline
& \textbf{Sistema} \\
\hline
5 & Consulta los datos según los filtros aplicados. \\
\hline
6 & Para reporte de banco o caja: calcula el saldo corriente en cada fila. \\
\hline
7 & Muestra los resultados en una tabla paginada. \\
\hline
& \textbf{Usuario} \\
\hline
8 & Presiona \textbf{Exportar PDF} o \textbf{Exportar Excel}. \\
\hline
& \textbf{Sistema} \\
\hline
9 & Genera el archivo en el formato seleccionado y lo descarga. \\
\hline
\end{tabularx}

\vspace{0.3cm}

\textbf{Precondiciones:} Deben existir datos registrados para el tipo de reporte seleccionado.

\textbf{Postcondición:} Los datos se muestran en pantalla. El archivo exportado se descarga al equipo del usuario.

\newpage

\subsubsection{Caso de Uso CU-15: Visualizar dashboard}

\begin{tabularx}{\textwidth}{|l|X|}
\hline
\rowcolor[gray]{0.9}
\textbf{Campo} & \textbf{Valor} \\
\hline
Caso de Uso & Visualizar dashboard \\
\hline
Actor & Gerente / Administrador \\
\hline
Descripción & El Gerente accede al dashboard ejecutivo que muestra indicadores clave (KPIs) del estado financiero de la empresa: total de proveedores, cooperativas, movimientos, retenciones del mes, saldo de caja, saldo bancario, gráficos de tendencias y distribución de retenciones por tipo. \\
\hline
\end{tabularx}

\vspace{0.3cm}

\textbf{Flujo de eventos básico}

\begin{tabularx}{\textwidth}{|c|X|}
\hline
& \textbf{Usuario} \\
\hline
1 & Inicia sesión y es redirigido al dashboard. \\
\hline
& \textbf{Sistema} \\
\hline
2 & Consulta los indicadores: total de proveedores, cooperativas, movimientos registrados. \\
\hline
3 & Calcula el total de retenciones del mes en curso y el monto acumulado. \\
\hline
4 & Consulta el saldo actual de caja y el saldo bancario consolidado. \\
\hline
5 & Genera los gráficos: tendencia de movimientos (12 meses), balance por cuenta, distribución de retenciones (S vs G), tendencia de retenciones (6 meses). \\
\hline
6 & Muestra las tarjetas KPI, tablas de movimientos recientes y retenciones recientes, y los gráficos. \\
\hline
\end{tabularx}

\vspace{0.3cm}

\textbf{Precondiciones:} Debe haber datos registrados en el sistema (proveedores, movimientos, retenciones).

\textbf{Postcondición:} Se visualiza el dashboard con información actualizada. No se modifica ningún dato.

\newpage

\subsection{Módulo de administración}

\subsubsection{Caso de Uso CU-16: Gestionar usuarios}

\begin{tabularx}{\textwidth}{|l|X|}
\hline
\rowcolor[gray]{0.9}
\textbf{Campo} & \textbf{Valor} \\
\hline
Caso de Uso & Gestionar usuarios \\
\hline
Actor & Administrador del Sistema \\
\hline
Descripción & El Administrador crea, modifica, consulta o elimina usuarios del sistema. Cada usuario tiene datos personales (nombre, apellido, CI, email, teléfono), puede tener uno o varios roles asignados, y puede estar activo o inactivo. \\
\hline
\end{tabularx}

\vspace{0.3cm}

\textbf{Flujo de eventos básico}

\begin{tabularx}{\textwidth}{|c|X|}
\hline
& \textbf{Usuario} \\
\hline
1 & Ingresa al módulo de usuarios. \\
\hline
2 & Presiona \textbf{Nuevo usuario}. \\
\hline
3 & Ingresa los datos del usuario: nombre, apellido, email, CI, teléfono, fecha de nacimiento, género. \\
\hline
4 & Asigna uno o varios roles al usuario. \\
\hline
5 & Presiona \textbf{Guardar}. \\
\hline
& \textbf{Sistema} \\
\hline
6 & Valida que el email y el CI no estén duplicados. \\
\hline
7 & Genera una contraseña temporal (por defecto el número de CI). \\
\hline
8 & Guarda el usuario con los roles asignados. \\
\hline
9 & Muestra mensaje de confirmación. \\
\hline
\end{tabularx}

\vspace{0.3cm}

\textbf{Precondiciones:} El Administrador debe tener permisos para gestionar usuarios.

\textbf{Postcondición:} El usuario queda registrado con sus roles. Puede iniciar sesión con su email y la contraseña generada.

\newpage

\subsubsection{Caso de Uso CU-17: Gestionar roles y permisos}

\begin{tabularx}{\textwidth}{|l|X|}
\hline
\rowcolor[gray]{0.9}
\textbf{Campo} & \textbf{Valor} \\
\hline
Caso de Uso & Gestionar roles y permisos \\
\hline
Actor & Administrador del Sistema \\
\hline
Descripción & El Administrador crea, modifica o elimina roles, y asigna o revoca permisos a cada rol. Los permisos definen qué funcionalidades del sistema puede acceder cada usuario según su rol. \\
\hline
\end{tabularx}

\vspace{0.3cm}

\textbf{Flujo de eventos básico}

\begin{tabularx}{\textwidth}{|c|X|}
\hline
& \textbf{Usuario} \\
\hline
1 & Ingresa al módulo de roles. \\
\hline
2 & Selecciona un rol existente o crea uno nuevo. \\
\hline
3 & Presiona \textbf{Asignar permisos}. \\
\hline
4 & Marca o desmarca los permisos deseados de la lista disponible. \\
\hline
5 & Presiona \textbf{Sincronizar}. \\
\hline
& \textbf{Sistema} \\
\hline
6 & Actualiza los permisos del rol seleccionado. \\
\hline
7 & Muestra mensaje de confirmación. \\
\hline
\end{tabularx}

\vspace{0.3cm}

\textbf{Precondiciones:} El Administrador debe tener permisos para asignar roles y permisos.

\textbf{Postcondición:} Los permisos del rol quedan actualizados. Los usuarios con ese rol ven reflejados los cambios en su próximo acceso.

\newpage

\section{Contraste con la situación actual}

\subsection{Validación de casos de uso}

\begin{longtable}{|L{4cm}|L{5cm}|L{5cm}|}
\hline
\rowcolor[gray]{0.9}
\textbf{Caso de uso} & \textbf{Contraste con entrevistas} & \textbf{Contraste con observación} \\
\hline
\endhead
CU-02: Gestionar retenciones & Resuelve la queja sobre desorganización y dispersión de registros al centralizar la información en una base de datos. & Elimina el paso de buscar archivos Excel en carpetas compartidas. \\
\hline
CU-03: Calcular descuentos automáticamente & Resuelve el problema de errores en cálculos manuales de retenciones mencionado por el Contador. & Automatiza el cálculo que antes se hacía con calculadora y Excel. \\
\hline
CU-05: Gestionar movimientos bancarios & Resuelve la queja sobre el 40\% del tiempo dedicado a tareas manuales en Excel. & Elimina el registro manual y el arrastre de fórmulas para calcular saldos. \\
\hline
CU-08: Gestionar movimientos de caja & Resuelve la dificultad de mantener registros exactos y actualizados. & Elimina las versiones contradictorias del Excel compartido por correo. \\
\hline
CU-11: Consultar estado de cuenta & Resuelve la necesidad de acceso rápido a saldos de proveedores y clientes. & Elimina la búsqueda manual en múltiples archivos. \\
\hline
CU-13: Generar reportes & Resuelve la necesidad de información actualizada para toma de decisiones. & Elimina la consolidación manual de datos para generar reportes mensuales. \\
\hline
CU-15: Visualizar dashboard & Resuelve la falta de visibilidad del estado financiero en tiempo real. & Reemplaza los informes esporádicos con indicadores actualizados automáticamente. \\
\hline
\end{longtable}

\subsection{Cumplimiento de objetivos específicos}

\begin{longtable}{|L{3cm}|L{5cm}|L{6cm}|}
\hline
\rowcolor[gray]{0.9}
\textbf{Objetivo específico} & \textbf{Casos de uso relacionados} & \textbf{Aporte} \\
\hline
Módulo de retenciones con cálculo automático & CU-02, CU-03, CU-04 & Automatiza el cálculo de descuentos y la emisión de recibos. \\
\hline
Módulo de libro de bancos con numeración fiscal & CU-05, CU-06, CU-07 & Centraliza los movimientos bancarios con numeración automática por año fiscal. \\
\hline
Módulo de caja con numeración mensual & CU-08, CU-09, CU-10 & Registra ingresos/egresos con numeración automática por mes. \\
\hline
Módulo de estados de cuenta & CU-11, CU-12 & Visualiza saldos e historial por cliente y proveedor. \\
\hline
Módulo de reportes PDF/Excel & CU-13, CU-14, CU-15 & Exporta información financiera para la toma de decisiones. \\
\hline
\end{longtable}

\end{document}
